\lyseis{Β3.7}{Σχετικές θέσεις ευθειών στο επίπεδο}
\sol{1} α) Σωστό. β) Λάθος: δύο τμήματα που δεν ακουμπούν μπορεί να βρίσκονται σε τεμνόμενες ευθείες, με σημείο τομής εκτός των τμημάτων. γ) Σωστό. δ) Σωστό (αίτημα του Ευκλείδη). ε) Σωστό.
\sol{2} α) παράλληλες \quad β) παράλληλες \quad γ) τεμνόμενες \quad δ) παράλληλες \quad ε) τεμνόμενες (σχηματίζουν γωνία $90^\circ$, άρα είναι μάλιστα κάθετες).
\sol{3} Ελεύθερη κατασκευή. Οι δύο ευθείες είναι παράλληλες προς την $\varepsilon$, άρα παράλληλες μεταξύ τους. Ταυτίζονται όταν τα $\mathrm{Α}$ και $\mathrm{Β}$ βρίσκονται πάνω στην ίδια παράλληλη προς την $\varepsilon$, δηλαδή όταν η $\mathrm{ΑΒ}$ είναι παράλληλη προς την $\varepsilon$.
\sol{4} Ελεύθερη κατασκευή. Η δεύτερη κάθετη είναι παράλληλη προς την $\varepsilon$, γιατί και οι δύο είναι κάθετες στην ίδια ευθεία.
\sol{5} Ελεύθερη κατασκευή. Η γωνία $\mathrm{ΑΚΒ}$ βγαίνει $50^\circ$, ίση με τη γωνία $x\hat{\mathrm{Ο}}y$.
\sol{6} Ελεύθερη κατασκευή. Υπάρχουν δύο ευθείες, μία σε κάθε μεριά της $\varepsilon$, στην απόσταση των $4$ cm.
\sol{7} Ελεύθερη κατασκευή. Αν οι ευθείες τους τέμνονται εκτός χαρτιού, τα τμήματα δεν είναι παράλληλα, παρόλο που δεν ακουμπούν.
\sol{8} Μεταξύ πέντε σειρών υπάρχουν $4$ διαστήματα των $40$ cm: $4\cdot40=160$ cm, δηλαδή $1{,}6$ μ.
\sol{9} Πάλι $1{,}5$ μ.: η απόσταση δύο παραλλήλων είναι η ίδια παντού, αλλιώς οι ράγες θα πλησίαζαν και θα συναντιούνταν.
\sol{10} Είναι παράλληλες, γιατί είναι κάθετες στην ίδια ευθεία.
\sol{11} Οι αρμοί της ίδιας κατεύθυνσης είναι παράλληλοι, γιατί είναι κάθετοι στους αρμούς της άλλης κατεύθυνσης (δηλαδή κάθετοι στην ίδια ευθεία).
\sol{12} Βάζουμε τον χάρακα πάνω στη δοσμένη ευθεία. Ακουμπάμε τον γνώμονα με μια πλευρά στον χάρακα. Σύρουμε τον γνώμονα κατά μήκος του χάρακα μέχρι η άλλη του πλευρά να περάσει από το σημείο, και χαράσσουμε. Η ευθεία που προκύπτει είναι παράλληλη προς τη δοσμένη.
\sol{13} Αν είχαν δύο κοινά σημεία, από τα δύο αυτά σημεία θα περνούσαν δύο διαφορετικές ευθείες, ενώ από δύο σημεία περνά μία μόνο ευθεία. Άρα έχουν το πολύ ένα κοινό σημείο.
\sol{14} Είναι παράλληλες μεταξύ τους. Αν είχαν σημείο τομής, από το σημείο αυτό θα περνούσαν δύο παράλληλες προς την $\varepsilon_2$, κάτι που αντιβαίνει στο αίτημα του Ευκλείδη.
\sol{15} Μία που δεν την τέμνει (η παράλληλη) και άπειρες που την τέμνουν.
\sol{16} Όχι, οι τεμνόμενες δεν είναι πάντα κάθετες (π.χ. δύο ευθείες που τέμνονται σε γωνία $48^\circ$). Ναι, οι κάθετες είναι πάντα τεμνόμενες (έχουν κοινό σημείο).
\sol{17} Ελεύθερη απάντηση. Π.χ. παράλληλες: οι άκρες του ραφιού και τα όρια ενός θρανίου, κάθετες: η γωνιά του τετραδίου.
\sol{18} Οι δύο κάθετες πτυχές της ίδιας πτυχής είναι παράλληλες μεταξύ τους, γιατί είναι κάθετες στην ίδια ευθεία.
